\documentclass[a4paper,twoside]{article}


% --- RUIMTE EN LAYOUT ---
\usepackage{setspace}

\usepackage{geometry} %size regulating
\geometry{
 a4paper,
 left=25mm,
 right=25mm,
 top=20mm,
 bottom=20mm,
 }

\usepackage{parskip} % for better spacing between paragraphs
\usepackage{flushend} % Alleen nodig als je laatste pagina van artikel wilt vullen met tekst, zodat de kolommen even lang zijn. Anders kan je dit beter weglaten.

% --- WISKUNDE ENSYMBOLEN ---
\usepackage{amsmath, amssymb, amsfonts, amsthm} % For nicer math format
\usepackage{mathrsfs} % For curly arrr
\usepackage{bm} % For bold symbols eg \bm{\nabla}
\usepackage{esint} % For bolletje in integrals
\usepackage{dsfont}
\usepackage{siunitx} % Voeg deze expliciet toe voor die graden en eenheden!
\usepackage{braket} % For bra-ket notation
\usepackage{mathtools}
\usepackage{url}
\allowdisplaybreaks

% --- OVERIGE ---

\usepackage{graphicx} % Required for inserting images
\usepackage[hidelinks]{hyperref} % geen kaders om linken
\usepackage{subcaption}
%\usepackage[dutch]{babel}
\usepackage[T1]{fontenc}
\usepackage[utf8]{inputenc}
\usepackage{lmodern}
\usepackage{multicol}


% --- Tikz packages ---
\usepackage{tikz}
\usetikzlibrary{angles,quotes, babel} % for pic (angle labels)
\usetikzlibrary{arrows.meta,decorations.markings,calc, decorations.pathmorphing, decorations.pathreplacing, positioning, shapes, shadings, patterns}
\usepackage{xcolor}
\usepackage{tikz-3dplot}
\usepackage[siunitx]{circuitikz}

\usepackage{etoolbox} % ifthen
\usepackage[outline]{contour} % glow around text
\usepackage{xfp} % higher precision (16 digits?)
\contourlength{1.1pt}



% --- THEOREM STIJLEN MET EXTRA RUIMTE ---

\theoremstyle{definition}
\newtheorem{postulate}{Postulaat}
\newtheorem{definition}{Definitie}
\newtheorem{derivation}{Afleiding}
\newtheorem{proposition}{Propositie}
\newtheorem{example}{Voorbeeld}

\theoremstyle{plain}
\newtheorem{theorem}{Theorem}


% --- EXTRA COMMANDO'S ---

\newcommand{\dbar}{{\mkern+3mu\mathchar'26\mkern-12mu d}}

% --- Titel en Introductie --- 

\title{Complexe Analyse: Moeilijkheden}
\author{T. Cools}
\date{2026-2027}


\begin{document}
\maketitle
\begin{abstract}
    In dit document hou ik bij met welke oefeningen in nog moeite heb, welke ik nog niet zo goed begrijp en duidelijk nog eens zal moeten naar kijken.
    Dus dit document is redelijk persoonlijk want iedereen ervaart andere moeilijkheden.
\end{abstract}

\section{Oefeningen Hoofdstuk 1}

\subsection{Algebra of Complex Numbers}

\subsection{Complex conjugates}

\subsection{Vectors}

\subsection{Polar coordinates}

Nog moeite met het begrip argument van een complex getal.
Begrijp mar half hoe dit werkt.

\subsection{The complex exponential}
\subsubsection{Oefening 1}
Moet ik mijn antwoord in \(\sin\) en \(\cos\) laten?
Of moet ik ze expliciet uitrekenen?
\subsubsection{Oefening 2}
Dit lukte in principe wel, maar ik zit fout met mijn kwadranten waardooor ik foutieve resultaten krijg.

Zeker nog eens naar kijken!

\subsubsection{Oefening 3}
Bewijs oke denk ik, nog eens controleren.

Mag ik gebruik maken van het feit dat \(e^{\pi i} = -1\), of moet ik dat aantonen?

\subsubsection{Oefening 4}

Idem vorig puntje.

\subsubsection{Oefening 5}

Nog te doen!

\subsubsection{Oefening 6}
Mijn bewijs nog controleren!

\subsubsection{Oefening 7}
(zie python plotjes)

\subsection{Powers and roots}
\subsubsection{Oefening 1}
Vind ik vaag wat ik moet doen\dots
\subsubsection{Oefening 2}
Nog te verbeteren
\subsubsection{Oefening 3}
WTF??
\subsection{Planar sets}
Oefeningen zijn te lang, skip ff.

\section{Oefeningen Hoofdstuk 2}
\subsection{Functions of a Complex Variable}
\subsubsection{Oefening 1}
Deelvraag c nog niet helemaal uitgerekend om tijd te besparen.
\subsubsection{Oefening 2}
Vervolg op oefening 1, ik wil eerst zeker zijn van de oplossing.
\subsubsection{Oefening 3}
\subsubsection{Oefening 4}
Zie python plots.
\subsubsection{Oefening 5}
\subsubsection{Oefening 6}
zie python plots.

\subsection{Limits and Continuity}
\subsubsection{Oefening 1}
\subsubsection{Oefening 2}
OK!
\subsubsection{Oefening 3}
nog verbeteren
\subsubsection{Oefening 4}

\subsubsection{Oefening 5}
\subsubsection{Oefening 6}
\subsubsection{Oefening 7}
\subsubsection{Oefening 8}
nog verbeteren
\subsection{Analyticity}
\subsubsection{Oefening 1}
HUH?

Na eerst de andere oefeningen te doen, is dit nu wel gelukt.

Algemeen: \(z = x+iy\)

door \(f(z) = Re(z) = x \implies v(x,y) = 0 \land u(x,y) = 1 \)
volgens cauchy-riemann moet dan \(1 = 0\), wat strijdig is.

We moeten de oef volgens onze def:
\[
f'(z_0) = \lim_{\Delta z \to 0} \frac{f(z_0 + \Delta z) - f(z_0)}{\Delta z}
\]
oplossen.

\[
= \lim_{\Delta z \to 0} \frac{x_0 + \Delta x - x_0}{\Delta z} 
= \lim_{\Delta z \to 0} \frac{\Delta x}{\Delta z}. 
\]
met \(\Delta z = \Delta x + i \Delta y\)
Dan 2 gevallen: \(\Delta x \to 0\) en \(\Delta_y = 0\)  en eens \( \Delta y \to 0 \) en \(\Delta x = 0\) bekijken.
\subsubsection{Oefening 2}
Lukt wel, nog verbeteren.
\subsubsection{Oefening 3}
\subsubsection{Oefening 4}
\subsection{The Cauchy-Riemann Equations}
\subsubsection{Oefening 1}
Ok! Wel goed verschil tussen differentiability en analyticity kennen!
\subsubsection{Oefening 2}
Gelukt!
\subsubsection{Oefening 3}
\subsubsection{Oefening 4}
Ok, maar vergeet niet:
\[
\boxed{f'(z_0) = \frac{\partial u(x_0, y_0)}{\partial x} + i\frac{\partial v(x_0, y_0)}{\partial x}} 
\qquad \text{ en } \qquad
 \boxed{f'(z_0) = -i\frac{\partial u(x_0, y_0)}{\partial y} + \frac{\partial v(x_0, y_0)}{\partial y}}
 \]
 m.a.w.:
 \[
 f'(z) = f'(u+iv) = u_x + i v_x = -i u_y + v_y.
 \]
\subsection{Harmonic Functions}
\subsubsection{Oefening 1}
Om te voldoen aan de laplace vergelijking moet gelden:
\[
\frac{\partial^2 u}{\partial x^2} + \frac{\partial^2 u}{\partial y^2} = 0
\qquad \text{ en } \qquad
\frac{\partial^2 v}{\partial x^2} + \frac{\partial^2 v}{\partial y^2} = 0
\]
m.a.w.:
\[
u_{xx} + u_{yy} = 0 \qquad \text{ en } \qquad v_{xx} + v_{yy} = 0  
\]
en belangrijk, \textbf{NIET} aan : \(u_xx + v_yy = 0\)! Dit heb ik perongelijk eerst gedaan.

\subsubsection{Oefening 2}
\subsubsection{Oefening 3}
Ik heb deze met een transformatie:

\(v \to -u'\) en \(u \to v'\) aangepakt om aaan te tonen dat:

\(u_x' = v_y'\) en \(u_y' = - v_x'\), wat de cauchy-riemann vergelijkingen zijn.
Dus \(u'\) en \(v'\) zijn harmonic conjugates.
\subsubsection{Oefening 4}
Yes die is gelukt! Wel te veel uitgewerkt, ik moest gewooon van het begin al enkel kijken naar \(v(x,y)\).
\subsubsection{Oefening 5}
\subsubsection{Oefening 6}

\end{document}